\documentclass[12pt,a4paper]{report}
\usepackage[]{amsmath,amsthm,amssymb,amscd}
\usepackage[latin1]{inputenc}
%\usepackage{fullpage}

\topmargin 0pt
\advance \topmargin by -\headheight
\advance \topmargin by -\headsep
     
\textheight 246mm
     
\oddsidemargin 0pt
\evensidemargin \oddsidemargin
\marginparwidth 0.5in
     
\textwidth 159mm


%               Theorem Environments
\theoremstyle{definition}
\newtheorem{ex}{}
\newcommand{\pd}[1]{\frac{\partial}{\partial #1}} %\pd{x}
%               Subquestions numbered with letters
\renewcommand{\theenumi}{\textbf{\alph{enumi}}}

%               Maths definitions

\newcommand{\NN}{\ensuremath{\mathbb N}}
\newcommand{\ZZ}{\ensuremath{\mathbb Z}}
\newcommand{\CC}{\ensuremath{\mathbb C}}
\newcommand{\TT}{\ensuremath{\mathbb T}}
\newcommand{\RR}{\ensuremath{\mathbb R}}
\newcommand{\QQ}{\ensuremath{\mathbb Q}}
\newcommand{\g}{\ensuremath{\frak{g}}}
\newcommand{\h}{\ensuremath{\frak{h}}}
\newcommand{\ft}{\ensuremath{\frak{t}}}
\newcommand{\cB}{\mathcal{B}}
\newcommand{\cN}{\mathcal{N}}
\newcommand{\cL}{\mathcal{L}}
\newcommand{\cD}{\mathcal{D}}
%\newcommand{\pd}[1]{\frac{\partial}{\partial #1}} %\pd{x}
%\newcommand{\pd}[1]{\frac{\partial}{\partial #1}} %\pd{x}

\begin{document}
\noindent
{\sffamily\large Marco Zambon 
\hfill     27/03/2012\\Geometr\'ia III, curso 2011-12}
 

%\noindent
%\hrulefill{}\\
\begin{center}\bfseries\large\textsf{Soluci\'on al examen Parcial    \vspace{-2mm}}
\end{center}

 
%\hrulefill{}\\

 \noindent
\hrulefill{}\\
 
 
\begin{ex}[8 puntos]    
Considera $S^2=\{v\in \RR^3: |v|=1\}$  con la relaci\'on de equivalencia $$v\sim_a w \Leftrightarrow (v=w \text{ o bien } v=-w).$$
Considera tambi\'en  $\RR^3-\{0\}$ con la relaci\'on de equivalencia $$v\sim_b w \Leftrightarrow \exists \lambda \neq 0: v=\lambda \cdot  w.$$
Encuentra \emph{explicitamente} un homeomorfismo $$S^2/\sim_a \to (\RR^3-\{0\})/\sim_b,$$ y demuestra que es un homeomorfismo.

\noindent{\bf Nota:} En clase vimos que existe un tal homeomorfismo
(ambos cocientes son homeomorfos al plano proyectivo). Aqu\'i os pido de escribir \emph{explicitamente} un tal homeomorfismo.
\end{ex}

\noindent{\bf Soluci\'on:} \begin{it}
Denota $\pi_a \colon S^2\to S^2/\sim_a$ la proyecci\'on canonica, y para $v\in S^2$ denota
$[v]_a:=\pi_a(v)$ (la clase de equivalencia de $v$). De manera similar, denota $\pi_b \colon \RR^3-\{0\}\to \RR^3-\{0\}/\sim_b$ la proyecci\'on canonica, y para $z\in \RR^3-\{0\}$ denota
$[z]_b:=\pi_b(z)$.
Considera
$$F \colon  S^2/\sim_a \to \RR^3-\{0\}/\sim_b,\;\; [v]_a \to [v]_b.$$ Esta aplicaci\'on est\'a bien definida, porque $[v]_b=[-v]_b$ para todo $v\in S^2$ (ya que $-v=(-1)\cdot v$). 

$F$ is injectiva: si para $v,w\in S^2$ existe  $\lambda \neq 0$ tal que $v=\lambda \cdot w,$ entonces $|v|=1=|w|$ implica que necesariamente $\lambda=1$ o $\lambda=-1$, y por lo tanto $[v]_a=[w]_a$.

 $F$ es sobreyectiva: dado $z\in \RR^3-\{0\}$, tenemos $F([\frac{z}{|z|}]_a)=[\frac{z}{|z|}]_b=[z]_b$. 
Entonces $F$ es biyectiva, y $F^{-1}([z]_b)=[\frac{z}{|z|}]_a$.


La aplicaci\'on $f\colon S^2 \to \RR^3-\{0\}, v\mapsto v$ es continua. $$\pi_b \circ f\colon  S^2 \to (\RR^3-\{0\})/\sim_b$$
satisface $(\pi_b \circ f)=F \circ \pi_a$. Como $\pi_b \circ f$ es continua, $F$ tambi\'en lo es (por el ejercicio 4 de la hoja 1).

La aplicaci\'on $g\colon \RR^3-\{0\} \to S^2, z\mapsto \frac{z}{|z|}$ es
continua. $$\pi_a \circ g \colon \RR^3-\{0\} \to S^2/\sim_a $$
satisface
satisface $(\pi_a \circ g)=F^{-1} \circ \pi_b$. Como $\pi_a \circ g$ es continua, $F^{-1}$ tambi\'en lo es (por el ejercicio 4 de la hoja 1).
En conclusi\'on, $F$ es un homeomorfismo.
\end{it}


\begin{ex} [7 puntos]
Considera $$S:=([0,1]\times[0,1])-\cup_{i=1}^5 D_i,$$ donde  $D_1,\dots,D_5$ son  discos abiertos en  $(0,1)\times(0,1)$ tal que  $\bar{D_i}\cap \bar{D_j}=\emptyset$ si $i\neq j$. Considera 
$S/\sim$, donde la relaci\'on de equivalencia $\sim$ es la siguiente:  cada punto es equivalente a s\'i mismo, y adem\'as $(0,y)\sim (1,1-y)$ para todo $y\in [0,1]$. 

> A qu\'e superficie en la lista de superficies compactas, conexas, con borde es homeomorfa  $S/\sim$? Demuestralo.
\end{ex}
\noindent{\bf Soluci\'on:} \begin{it}
$S/\sim$ es homeomorfo a $$P_2(\RR)- \cup_{i=1}^6 D_i',$$ donde
los $D_i'$ son  discos abiertos en $P_2(\RR)$ tal que sus cierres son disjuntos dos a dos.

Est\'a claro que $S$ es la banda de M\"obius  $M$ menos 5 discos abiertos. Ahora utilizamos el hecho, visto en clase, que pegando la banda de M\"obius $M$ y un disco cerrado a lo largo de sus bordes se obtiene $P_2(\RR)$. Entonces $M\cong P_2(\RR)- D_1'$, donde  $D_1'$ es un disco abierto en $P_2(\RR)$. Por lo tanto $S$ es homeomorfo a $P_2(\RR)$ menos 6 discos abiertos.
\end{it} 

\begin{ex}  [8 puntos]
Sea $S^2=\{v\in \RR^3: |v|=1\}$, con su estructura   diferencial standard. Sea $\textbf{S}:=(0,0,-1)\in S^2$.
Para todo $(x,y,z)\in \RR^3$, se denota por $|(x,y,z)|:=\sqrt{x^2+y^2+z^2}$ su norma.

> Es la funci\'on
$$f \colon S^2\to \RR,  \;\;\; v\mapsto |v-\textbf{S}|$$
diferenciable? Demuestralo.

\noindent{\bf Nota:} Puedes utilizar que una carta diferencial de $S^2$ es dada por la proyecci\'on estereogr\'afica $\phi_N$, y que su inversa es $$(\phi_N)^{-1} \colon \RR^2 \to S^2-\{\textbf{N}\},\;\;(u,v)\mapsto \frac{1}{u^2+v^2+1}(2u,2v,u^2+v^2-1),$$
donde $\textbf{N}:=(0,0,1)$.
\end{ex}
\noindent{\bf Soluci\'on:} \begin{it}
No, $f$ no es diferenciable en el punto $\textbf{S}$.

Primero, escribimos $f$ as\'i: 
$$f((x,y,z))=|(x,y,z)-\textbf{S}|=\sqrt{x^2+y^2+(z+1)^2}=
\sqrt{1+2z+1}=\sqrt{2(z+1)}.$$
Tomamos la carta $\phi_N$. Obtenemos
$$f \circ (\phi_N)^{-1} \colon \RR^2 \to \RR, (u,v)\mapsto \sqrt{2(\frac{u^2+v^2-1}{u^2+v^2+1}+1)}=2\sqrt{\frac{u^2+v^2}{u^2+v^2+1}}
,$$ es decir, $(f \circ (\phi_N)^{-1})(X)=2\frac{|X|}{\sqrt{|X|^2+1}}$ para todo $X\in \RR^2$.
Esta funci\'on no es diferenciable en el origen de $\RR^2$ (ya que la norma no lo es).
\end{it}




\begin{ex}[7 puntos]
Si $M$ es una variedad diferencial y $S,R$ son subvariedades de $M$ tal que $dim(S)=dim(R)$, > es verdad que la union $S\cup R$ es siempre  una subvariedad de $M$? Demuestralo.
\end{ex}
\noindent{\bf Soluci\'on:} \begin{it}
No.

Un contraej\'emplo es lo siguiente. Sea $M=\RR^2$ (con su estructura standard de variedad diferencial). Sea $S=\RR\times \{0\}$ (el eje $x$), y $R=\{0\}\times \RR$ (el eje $y$). $S$ y $R$ son subvariedades de $M$ de dimensi\'on $1$: $(\RR^2,Id_{\RR^2})$ es una carta de $M$ adaptada a $S$, y $(\RR^2,(x,y)\mapsto(y,x))$ es una carta de $M$ adaptada a $R$.

La uni\'on $S\cup R$ no es una subvariedad de $M$. Una manera de verlo es: si lo fuera, entonces $S\cup R$ ser\'ia una variedad diferencial de dimensi\'on $1$.
Pero existe un entorno conexo $U$ del origen $0$ en $S\cup R$  tal que $U-\{0\}$ tiene $4$ componentes conexas, entonces $U$ no puede ser homeomorfo a un intervalo.
\end{it}
 

%\begin{ex}  
%Demuestra que 
%$$T \;\;\sharp \;\; C\cong T-(D_1 \cup D_2)$$
%donde $T=S^1\times S^1$ es el toro, $C=S^1\times [0,1]$ el cilindro, y $D_i\subset T$ (por $i=1,2$) son discos cerrados tal que $D_1\cap D_2=\emptyset$.
%\end{ex} 






 \end{document}
 %%%%%%%
 \begin{ex}

\end{ex}

 \begin{ex}
\begin{enumerate}
\item 
\end{enumerate}
\end{ex}
 
 
 
 